\renewcommand{\bibname}{Quellen und Bandzuordnung}
\begin{thebibliography}{9}
\addcontentsline{toc}{chapter}{Quellen und Bandzuordnung}
\footnotesize
\setlength{\itemsep}{5pt}
\setlength{\parskip}{0pt}
\interlinepenalty=10000

\bibitem{project}
Martin Kunze: \emph{Die Grundlagen der Mathematik}.
Band~12: Halbordnungsaxiome, Dualität, Extremalelemente,
Restriktionen und Hauptfilter. Der Abschnitt über Präordnungen
und Quotienten wird nicht übernommen.
Mengenoperationen und endliche Folgen dienen als vorausgesetzte
Notation. Format nach den Leitfäden zu Band~3, Band~4, Band~5,
Band~28 und Band~37.

\bibitem{ker}
Andrew D. Ker: \emph{Discrete Mathematics}, University of Oxford,
Michaelmas Term 2010. Kapitel~8.1--8.2, Druckseiten~97--100:
Halbordnungen, strenge Ordnungen, Antiketten und Produktordnungen.
Vergleichsstellen zu Definitionen~\ref{b12guide:def:order},
\ref{b12guide:def:strict}, \ref{b12guide:def:antichain}
und Theorem~\ref{b12guide:thm:product}.
\url{https://www.cs.ox.ac.uk/andrew.ker/docs/discretemaths-lecture-notes-mt2010.pdf}.

\bibitem{minimal}
The mathlib Community: \emph{Mathlib.Order.Minimal},
offizielle Dokumentation, abgerufen am 23.~Oktober 2026.
Einträge \path{minimal_iff_forall_lt}, \path{Minimal.mono},
\path{IsLeast.minimal_iff} und \path{IsGreatest.maximal_iff}:
strenge Kriterien, Vererbung sowie Verhältnis zwischen Minimum
und minimalen Elementen in Kapitel~\ref{b12guide:ch:extremal}.
\url{https://leanprover-community.github.io/mathlib4_docs/Mathlib/Order/Minimal.html}.

\bibitem{upperlower}
The mathlib Community: \emph{Mathlib.Order.UpperLower.Basic},
offizielle Dokumentation, abgerufen am 23.~Oktober 2026.
Untere und obere Mengen, Komplement, Vereinigungen und Durchschnitte;
insbesondere \path{isLowerSet_iff_Iic_subset}.
Vergleichsstellen zu Theoremen~\ref{b12guide:thm:closed-criteria}
und~\ref{b12guide:thm:set-operations}.
\url{https://leanprover-community.github.io/mathlib4_docs/Mathlib/Order/UpperLower/Basic.html}.

\bibitem{closure}
The mathlib Community: \emph{Mathlib.Order.UpperLower.Closure},
offizielle Dokumentation, abgerufen am 23.~Oktober 2026.
Einträge \path{mem_upperClosure}, \path{mem_lowerClosure},
\path{upperClosure_min}, \path{upperClosure_sUnion},
\path{ordConnected_iff_upperClosure_inter_lowerClosure} und
\path{IsAntichain.minimal_mem_upperClosure_iff_mem}.
Vergleichsstellen für Hüllen, Konvexität und Antiketten in
Kapiteln~\ref{b12guide:ch:order-sets} und~\ref{b12guide:ch:subsets}.
Hier bezeichnet \glqq kleinste\grqq{} stets die gewöhnliche
Mengeninklusion, unabhängig von den Ordnungen auf Mengentypen
in der Bibliothek.
\url{https://leanprover-community.github.io/mathlib4_docs/Mathlib/Order/UpperLower/Closure.html}.

\end{thebibliography}

\begingroup
\footnotesize
\noindent\textbf{Einordnung.}
Die übrigen Kriterien, Rechengesetze und Konstruktionen sind
hier zusammenhängend formulierte elementare Folgerungen aus
den Halbordnungsaxiomen und den angegebenen Mengendefinitionen.
Dies betrifft insbesondere die Erzeugung durch endliche
Vergleichsfolgen und die explizite Erweiterung um einen Vergleich
in Kapitel~\ref{b12guide:ch:generated}.
Die Quellen dienen dem Abgleich von Begriffen und zentralen
Resultaten; nicht jede Formulierung steht dort wörtlich.
Es werden keine neuen Forschungsergebnisse und keine historischen
Erstentdeckungen beansprucht.
\endgroup
